% VI/41 reader bridge -- additive mathematical exposition.
\paragraph{Linear equivalence and divisor classes.}
For a Noetherian normal integral scheme $X$, let
$\operatorname{Div}(X)$ be the free abelian group of Weil divisors.
The principal divisors $\operatorname{div}(f)$, for
$f\in K(X)^\times$, form a subgroup.  Its quotient is the
\emph{divisor class group}
\[
 \operatorname{Cl}(X)
  =\operatorname{Div}(X)/\{\operatorname{div}(f):f\in K(X)^\times\}.
\]
Two Weil divisors $D,E$ are linearly equivalent precisely when
$D-E=\operatorname{div}(f)$ for some nonzero rational function.
Linear equivalence is \emph{not} equality of supports: adding a
principal divisor changes the zeros and poles while preserving the
class.  On a normal scheme Cartier divisor classes map naturally
to $\operatorname{Cl}(X)$, but this map need not be surjective at
singular points.  If $X$ is regular and Noetherian, the two class
theories agree and
$\operatorname{Pic}(X)\simeq\operatorname{Cl}(X)$.

\paragraph{Worked example 1: affine space has trivial class group.}
Let $A=k[x_1,\ldots,x_n]$.  It is a unique factorization domain.
Every codimension-one prime corresponds to a principal ideal $(p)$
for an irreducible polynomial $p$.  The corresponding prime divisor
is exactly $\operatorname{div}(p)$: it has coefficient one along
$V(p)$ and zero along the other height-one primes.  Therefore the
generators of $\operatorname{Div}(\operatorname{Spec}A)$ are all
principal, and
\[
 \operatorname{Cl}(\mathbb A_k^n)=0.
\]
This is a statement about \emph{Weil divisors modulo rational
functions}; it does not assert that there are no hypersurfaces in
affine space.  It says that every such divisor is equivalent to zero.

\paragraph{Worked example 2: the complete class-group calculation on
$\mathbb P^1$.}
For a closed point $P$ of $\mathbb P_k^1$, define its degree by
$\deg(P)=[\kappa(P):k]$.  Extend additively to Weil divisors:
$\deg(\sum n_P[P])=\sum n_P\deg(P)$.
If $f\in k(t)^\times$, the total degree of its zeros is equal to
the total degree of its poles, counted with the residue-field
multiplicities.  Thus $\deg\operatorname{div}(f)=0$, and degree
descends to a map $\operatorname{Cl}(\mathbb P_k^1)\to\mathbb Z$.

To see that it is an isomorphism, choose
$P_\infty=[1:0]$, of degree one.  Every closed point in the
chart $t=X/Y$ is represented by a monic irreducible
$p(t)\in k[t]$ of degree $d$.  On the projective line,
\[
 \operatorname{div}(p(t))=[P]-d[P_\infty].
\]
The coefficient at $P$ is one because $p$ is a uniformizer there;
the pole at infinity has order $d$.  Hence
$[P]\sim d[P_\infty]$.  Every divisor class is therefore an
integer multiple of $[P_\infty]$.  Because the degree map sends
$[P_\infty]$ to one, there is no nontrivial relation.  We obtain
\[
 \operatorname{Cl}(\mathbb P_k^1)\simeq\mathbb Z,
 \qquad[D]\longmapsto\deg D.
\]
For example, $[0:1]$ and $[1:0]$ are not equal as divisors, but
they are linearly equivalent by the rational function $t=X/Y$.

\paragraph{Worked example 3: hyperplane classes on $\mathbb P^2$.}
Write $H=V_+(X)$ and let $F(X,Y,Z)$ be a nonzero homogeneous
polynomial of degree $d$.  The fraction $F/X^d$ has homogeneous
degree zero and hence defines a rational function on
$\mathbb P_k^2$.  Its divisor is
\[
 \operatorname{div}(F/X^d)=[V_+(F)]-dH,
\]
where reducible or repeated factors of $F$ are counted with their
multiplicities.  Therefore every plane curve of degree $d$ is
linearly equivalent to $dH$.  In fact
$\operatorname{Cl}(\mathbb P_k^2)\simeq\mathbb Z$ with generator
$[H]$: the degree of a homogeneous hypersurface gives the
integer coefficient.  This is the projective analogue of the
factorization calculation on $\mathbb A_k^2$, but here the
hyperplane class does \emph{not} become principal.

\paragraph{Worked example 4: a singular surface with two-torsion.}
Return to $A=k[u,v,w]/(uv-w^2)$ with $k$ algebraically closed
and of characteristic different from $2$.
Let $D=V(u,w)$ be the height-one prime divisor from VI/40.
The hypersurface $V(u)$ has set-theoretic support $D$, since
$A/(u)\simeq k[v,w]/(w^2)$ has only one minimal prime $(w)$.
At the generic point of $D$, the element $v$ is a unit and
$u=w^2/v$, so the discrete valuation has
$\operatorname{ord}_D(u)=2$.  Consequently
\[
 \operatorname{div}(u)=2D,
 \qquad 2[D]=0\quad\text{in }\operatorname{Cl}(X).
\]
But VI/40 showed that the ideal $(u,w)$ is not locally principal
at the vertex.  If $D$ were principal as a Weil divisor, it would
be Cartier; therefore $[D]\ne0$.  Thus $[D]$ has \emph{exact}
order two.

For completeness, the localization $A_u$ is the Laurent
polynomial ring $k[u,u^{-1},w]$, because $v=w^2/u$ once $u$ is
invertible.  It is a UFD and has class group zero.  The localization
sequence for class groups shows that $\operatorname{Cl}(A)$ is
generated by height-one prime divisors supported in $V(u)$.
The only such prime divisor is $D$.  Combining this with the
previous paragraph gives the full answer
\[
 \operatorname{Cl}\bigl(k[u,v,w]/(uv-w^2)\bigr)
     \simeq\mathbb Z/2\mathbb Z.
\]
The nontrivial class is a genuine obstruction to expressing the
Weil divisor $D$ by a local equation at the singular vertex.

\paragraph{What to carry into the next sections.}
Weil divisors record codimension-one cycles; Cartier divisors
record local nonzerodivisor equations; line bundles record the
transition functions of these local equations.  On regular schemes
the viewpoints agree, but on singular schemes they can differ.
The class group remembers divisors up to rational functions, while
the Picard group remembers invertible sheaves.  Later calculations
with line bundles and intersections should state which group is
being used before assigning a numerical degree or an intersection
multiplicity.
